% \newglossaryentry{sma}{
%     name=Système Multi-Agents,
%     description={Système composé d'entités autonomes appelées agents qui interagissent entre elles et avec leur environnement}
% }

% \newglossaryentry{passi}{
%     name=PASSI,
%     description={Process for Agent Societies Specification and Implementation. Méthodologie de conception de systèmes multi-agents qui guide le développement depuis les exigences jusqu'au code}
% }

% \newglossaryentry{agent}{
%     name=Agent,
%     description={Entité logicielle autonome capable de percevoir son environnement, de raisonner et d'agir en fonction de ses objectifs}
% }

% \newglossaryentry{ontologie}{
%     name=Ontologie,
%     description={Représentation formelle des connaissances et des relations entre concepts dans un domaine spécifique}
% }

% \newglossaryentry{protocole}{
%     name=Protocole d'interaction,
%     description={Ensemble de règles définissant la séquence et le format des messages échangés entre agents}
% }

% \newglossaryentry{marche-financier}{
%     name=Marché financier,
%     description={Système qui permet l'échange d'instruments financiers comme les actions, les obligations et les produits dérivés}
% }

% \newglossaryentry{trader}{
%     name=Trader,
%     description={Acteur du marché financier qui achète et vend des instruments financiers pour générer un profit}
% }

% \newglossaryentry{courtier}{
%     name=Courtier (Broker),
%     description={Intermédiaire qui exécute les ordres d'achat et de vente pour le compte de ses clients sur les marchés financiers}
% }

% \newglossaryentry{regulateur}{
%     name=Régulateur,
%     description={Autorité chargée de surveiller les marchés financiers et d'assurer le respect des règles et réglementations}
% }

% \newglossaryentry{publish-subscribe}{
%     name=Modèle publish-subscribe,
%     description={Patron de conception où les expéditeurs de messages (publishers) ne programment pas les messages pour être envoyés directement à des destinataires spécifiques, mais catégorisent les messages publiés, et les récepteurs (subscribers) s'abonnent aux catégories qui les intéressent}
% }

